\chapter{Rappresentazione a grafo e schema delle feature}
\label{app:graph-features}

Questa appendice specifica la rappresentazione da scacchiera a grafo usata dai modelli neurali di GraphMate.  È una specifica implementativa, non una spiegazione generale delle reti neurali su grafi.  L'Appendice~\ref{app:hyperparameters} indica quali modelli consumano ciascuna rappresentazione, mentre l'Appendice~\ref{app:data-pipeline} descrive come i file di grafo sono prodotti a partire dai dati sorgente.

\section{Indicizzazione delle case e layout dei tensori}
\label{app:graph-features:indexing}

Ogni posizione è rappresentata come un grafo PyTorch Geometric con un nodo per ciascuna casa della scacchiera.  I nodi sono ordinati secondo l'indice di casa di python-chess, da \(a1=0\) fino a \(h8=63\).  La coordinata di riga è l'indice di traversa diviso per 7 e la coordinata di colonna è l'indice di colonna diviso per 7:
\[
r(s)=\frac{\operatorname{rank}(s)}{7}, \qquad c(s)=\frac{\operatorname{file}(s)}{7}.
\]
Pertanto la prima traversa ha coordinata di riga 0, l'ottava traversa ha coordinata di riga 1, la colonna \(a\) ha coordinata di colonna 0 e la colonna \(h\) ha coordinata di colonna 1.  L'oggetto dati del grafo contiene la matrice dei nodi \(x\in\mathbb{R}^{64\times15}\), l'indice degli archi sparse COO \(\texttt{edge\_index}\in\mathbb{N}^{2\times E}\), la matrice degli archi \(\texttt{edge\_attr}\in\mathbb{R}^{E\times5}\) e le feature globali \(\texttt{global\_features}\in\mathbb{R}^{4}\).  Il dataset conserva inoltre campi di metadati, tra cui FEN, mossa target, identificatore del puzzle, rating, profondità di matto e indice della riga sorgente.

\begin{figure}[H]
\centering
\fbox{\begin{minipage}{0.92\linewidth}
\centering
\texttt{Posizione FEN} \(\rightarrow\) \texttt{64 nodi-casa} \(\rightarrow\) \texttt{archi relazionali diretti} \(\rightarrow\) \texttt{grafo PyG} \(\rightarrow\) \texttt{contesto globale / candidate legali}
\end{minipage}}
\caption{Costruzione del grafo a livello concettuale.  I canali esatti di nodi, archi e feature globali sono specificati nelle Tabelle~\ref{tab:node-features}, \ref{tab:edge-features} e \ref{tab:global-features}.}
\label{fig:graph-pipeline}
\end{figure}


\section{Feature dei nodi}
\label{app:graph-features:nodes}

La Tabella~\ref{tab:node-features} riporta l'ordine esatto delle 15 dimensioni dei nodi.  Le case vuote ricevono indicatori di tipo pezzo pari a zero, colore 0, occupazione 0, pin 0 e valore del pezzo 0, mantenendo però le coordinate normalizzate e le feature di conteggio degli attacchi.  Il campo colore vale 1 per un pezzo bianco e 0 per un pezzo nero o per una casa vuota.  La feature di mobilità legale conta le mosse legali che hanno origine dalla casa dopo aver valutato copie della scacchiera con il tratto impostato prima al Bianco e poi al Nero; è quindi un conteggio di mobilità per casa di origine, non un'etichetta di mossa legale per il target di addestramento.  La feature di pezzo inchiodato è locale a un nodo occupato e non deve essere confusa con l'indicatore globale di scacco nella Tabella~\ref{tab:global-features}.

\begin{table}[H]
\centering
\scriptsize
\caption{Ordine delle feature dei nodi per ciascuno dei 64 nodi-casa.}
\label{tab:node-features}
\begin{tabular}{p{0.10\linewidth}p{0.23\linewidth}p{0.55\linewidth}}
\toprule
Indice & Feature & Codifica \\
\midrule
0--5 & Tipo di pezzo & One-hot per pedone, cavallo, alfiere, torre, donna, re; tutti zero per case vuote. \\
6 & Colore & 1 per pezzo bianco; 0 per pezzo nero o casa vuota. \\
7 & Occupata & 1 se una casa contiene un pezzo; altrimenti 0. \\
8 & Riga & \(\operatorname{rank}(s)/7\), con rank 0 sulla prima traversa. \\
9 & Colonna & \(\operatorname{file}(s)/7\), con file 0 sulla colonna \(a\). \\
10 & Attaccata dal Bianco & Conteggio degli attaccanti bianchi della casa nello stato corrente. \\
11 & Attaccata dal Nero & Conteggio degli attaccanti neri della casa nello stato corrente. \\
12 & Mobilità legale & Conteggio delle mosse legali originate dalla casa su stati copiati con tratto al Bianco e al Nero. \\
13 & Pezzo inchiodato & 1 se il pezzo occupante è inchiodato al proprio re; 0 per pezzi non inchiodati o case vuote. \\
14 & Valore del pezzo & Pedone 1, cavallo/alfiere 3, torre 5, donna 9, re 0, vuota 0. \\
\bottomrule
\end{tabular}
\end{table}


\section{Feature degli archi}
\label{app:graph-features:edges}

Gli archi sono diretti e multi-etichetta.  Se più relazioni si applicano alla stessa coppia ordinata di case, i cinque canali dell'arco sono aggregati tramite massimo componente per componente.  La lista risultante degli archi è memorizzata in formato sparse COO.  Gli archi di mossa legale puntano dall'origine alla destinazione per ogni mossa legale nella posizione.  Gli archi di attacco e difesa puntano da una casa sorgente occupata a una casa bersaglio occupata attaccata dal pezzo sorgente; il canale è attacco per colori opposti e difesa per lo stesso colore.  Gli archi di inchiodatura sono generati lungo le linee del re dal pezzo scorrevole attaccante al pezzo inchiodato.  Gli archi di linea di scacco sono generati solo quando il lato al tratto è sotto scacco e collegano l'attaccante che dà scacco al re; non enumerano tutte le case sulla linea geometrica.

\begin{table}[H]
\centering
\small
\caption{Canali delle feature degli archi diretti.  Relazioni multiple sulla stessa coppia diretta sono combinate tramite massimo componente per componente.}
\label{tab:edge-features}
\begin{tabular}{p{0.12\linewidth}p{0.22\linewidth}p{0.52\linewidth}}
\toprule
Indice & Canale & Semantica \\
\midrule
0 & Mossa legale & Arco origine-destinazione per ciascuna mossa legale nella posizione. \\
1 & Attacco & La sorgente occupata attacca un bersaglio occupato di colore opposto. \\
2 & Difesa & La sorgente occupata attacca un bersaglio occupato dello stesso colore. \\
3 & Inchiodatura & Pezzo scorrevole attaccante verso pezzo inchiodato lungo la linea del re. \\
4 & Linea di scacco & Attaccante che dà scacco verso re quando il lato al tratto è sotto scacco. \\
\bottomrule
\end{tabular}
\end{table}


\section{Feature globali}
\label{app:graph-features:globals}

Le quattro feature globali nella Tabella~\ref{tab:global-features} codificano lato al tratto, stato di scacco e contatori di mossa con clipping.  Le normalizzazioni dei contatori sono
\[
\text{fullmove\_norm}=\frac{\min(\text{fullmove},200)}{200}, \qquad
\text{halfmove\_norm}=\frac{\min(\text{halfmove},100)}{100}.
\]
Queste feature sono concatenate all'embedding del grafo ottenuto con pooling e sono conservate anche nei grafi successivi alla mossa usati da Model A4.

\begin{table}[H]
\centering
\small
\caption{Ordine delle feature globali del grafo.}
\label{tab:global-features}
\begin{tabular}{p{0.12\linewidth}p{0.28\linewidth}p{0.48\linewidth}}
\toprule
Indice & Feature & Codifica \\
\midrule
0 & Lato al tratto & 1 per Bianco al tratto, 0 per Nero al tratto. \\
1 & Sotto scacco & 1 se il lato al tratto è sotto scacco; altrimenti 0. \\
2 & Fullmove & \(\min(\text{fullmove},200)/200\). \\
3 & Halfmove & \(\min(\text{halfmove},100)/100\). \\
\bottomrule
\end{tabular}
\end{table}


\section{Rappresentazioni delle mosse candidate}
\label{app:graph-features:candidates}

I modelli basati su mosse candidate legali costruiscono un vettore candidato per ciascuna mossa legale.  Il raggruppamento delle candidate è memorizzato con un vettore di puntatori, così che la loss sia calcolata sulle mosse legali appartenenti alla stessa posizione.  Le promozioni usano un identificatore a cinque valori, con nessuna promozione mappata a 0 e donna, torre, alfiere e cavallo mappati rispettivamente a 1--4 prima dell'applicazione di un embedding di promozione a 8 dimensioni.

Model A3 concatena embedding della casa sorgente, embedding della casa destinazione, contesto del grafo ed embedding della promozione.  Model B aggiunge un embedding di timing al contesto del grafo.  Model A4 chiede prima al Model A3 congelato le Top-5 candidate, poi applica ciascuna candidata per produrre un grafo della posizione risultante, codifica tale grafo e concatena il contesto successivo alla mossa con la rappresentazione candidata congelata di A3 e due feature dello score A3.  La Tabella~\ref{tab:candidate-dimensions} riporta l'aritmetica verificata delle dimensioni.

\begin{table}[H]
\centering
\small
\caption{Dimensioni dei vettori candidati per i modelli a mosse legali e di reranking.}
\label{tab:candidate-dimensions}
\begin{tabular}{p{0.14\linewidth}p{0.58\linewidth}p{0.14\linewidth}}
\toprule
Modello & Componenti & Dimensione \\
\midrule
A3 & nodo sorgente 128 + nodo destinazione 128 + contesto grafo 132 + embedding promozione 8 & 396 \\
B & nodo sorgente 128 + nodo destinazione 128 + contesto grafo con timing \(132+16\) + embedding promozione 8 & 412 \\
A4 & candidato A3 congelato 396 + contesto post-mossa 132 + feature score A3 2 & 530 \\
\bottomrule
\end{tabular}
\end{table}

